% Source: tex/chapters/ch06/07_ソフトウェアエンジニアリング運用制御.tex
\subsubsection{変更管理}

変更管理とは、すべての変更を評価、承認、実装、レビューされた状態で行うためのプロセスである。変更要求は記録され、緊急、至急、重大、軽微などに分類される。変更管理では、変更のリスク、ロールバック戦略の必要性、関係者への通知、変更スケジュールを扱う。

従来型の開発では、多くの変更をまとめて定期リリースとして提供することが多かった。DevOpsでは、小さな変更単位を独立に、必要に応じて届けることを目指す。そのためには、ソフトウェアが小さく独立した配備に向くように設計されている必要がある。
